% ============================================================================
% Fichier  : partie2/P02_C05_osint.tex
% Titre    : Reconnaissance Passive et OSINT
% Auteur   : MINKA MI NGUIDJOÏ Thierry Emmanuel
% Version  : 1.0 -- novembre 2026
% ============================================================================
% Sources : PTES (collecte de renseignements) ; NIST SP 800-115 ; RFC 6962
% (Certificate Transparency) ; MITRE ATT&CK (tactique Reconnaissance).
% Résultats de l'Opération LOGONE : cas fictif.
% ============================================================================

\chapter{Reconnaissance Passive et OSINT}
\label{chap:osint}
\chaptermark{Reconnaissance passive et OSINT}

\epigraph{%
    « L'attaquant connaît souvent votre système d'information\\
    mieux que votre inventaire. »%
}{Principe du cours}

\niveauChapitre{Intermédiaire}{Chapitres~\ref{chap:engagement}
et~\ref{chap:methodes}.}

\begin{boxobjectifs}
\begin{itemize}
    \item Distinguer reconnaissance passive et active.
    \item Exploiter les sources ouvertes~: DNS, certificats publics,
          annonces, réseaux professionnels, dépôts de code.
    \item Construire la surface d'attaque externe d'une organisation.
    \item Rester dans le cadre légal et éthique de la collecte.
    \item Transformer les résultats en mesures de réduction de l'exposition.
\end{itemize}
\end{boxobjectifs}

\bigskip

% ============================================================================
\section*{Question directrice : que voyait-on de LiZbiZ depuis Internet, sans rien toucher~?}
% ============================================================================

Les 19 et 20~octobre, M.~TAMO ne se connecte à aucun système de LiZbiZ. Il
lit ce que le monde peut lire~: les enregistrements DNS, les journaux
publics de certificats, les offres d'emploi, les profils professionnels des
salariés. À la fin du deuxième jour, il connaît trois serveurs exposés que
l'inventaire de M.~KOUNOU ignore.

\medskip
\textbf{Plan du chapitre.}
\begin{enumerate}
    \item Passive ou active (section~\ref{sec:c05-passive}).
    \item Les sources ouvertes (section~\ref{sec:c05-sources}).
    \item La surface d'attaque externe (section~\ref{sec:c05-surface}).
    \item Les limites de la collecte (section~\ref{sec:c05-limites}).
    \item Réduire l'exposition (section~\ref{sec:c05-reduire}).
\end{enumerate}

\separateur

% ============================================================================
\section{Passive ou Active}
\label{sec:c05-passive}
% ============================================================================

\begin{boxdef}
\termecle{Reconnaissance passive}~: collecte d'informations sans interagir
avec les systèmes de la cible (sources publiques, services tiers).
\termecle{Reconnaissance active}~: collecte par interaction directe
(résolution de noms auprès des serveurs de la cible, balayage). Le
renseignement en sources ouvertes, ou \termecle{OSINT}, rassemble les
informations publiquement accessibles.
\end{boxdef}

La reconnaissance passive ne laisse presque aucune trace chez la cible~:
elle est aussi ce que fait un attaquant avant d'attaquer. ATT\&CK lui
consacre une tactique entière, la \emph{reconnaissance}~\cite{mitreattack}.

% ============================================================================
\section{Les Sources Ouvertes}
\label{sec:c05-sources}
% ============================================================================

\begin{tableaucours}{L{3.6cm}L{4.2cm}L{4.4cm}}{%
    \tableauHeader{Source} &
    \tableauHeader{Ce qu'elle révèle} &
    \tableauHeader{Exemple chez LiZbiZ}
}
Enregistrements DNS & Serveurs, messagerie, services & Serveurs de
messagerie hébergés en nuage \\
Journaux de transparence des certificats & Noms de domaine pour lesquels un
certificat a été émis & 14~sous-domaines de lizbiz.cm \\
Registres d'adresses & Plages attribuées, hébergeur & La plage
198.51.100.0/27 \\
Offres d'emploi & Technologies utilisées & « Administrateur de notre
pare-feu et de notre passerelle d'accès distant » \\
Réseaux professionnels & Organigramme, format des adresses &
prenom.nom@lizbiz.cm \\
Dépôts de code publics & Configurations, secrets oubliés & Aucun secret
trouvé \\
\end{tableaucours}
\vspace{-0.8em}
\begin{center}{\small\itshape Tableau~5.1. Sources ouvertes et apports}\end{center}

Les journaux de \termecle{transparence des certificats}, définis par la RFC~6962,
enregistrent publiquement les certificats émis par les autorités
participantes~\cite{rfc6962}. Ils servent à détecter les certificats émis
par erreur~; ils révèlent aussi les noms de serveurs qu'une organisation
croyait discrets.

% ============================================================================
\section{La Surface d'Attaque Externe}
\label{sec:c05-surface}
% ============================================================================

\begin{boxscenario}{Opération LOGONE. Ce que l'inventaire ignorait}
\textbf{Constatations} (20~octobre). Sur 14~sous-domaines trouvés dans les
journaux de certificats, 11 figurent à l'inventaire. Trois n'y sont pas~:
« ancien-portail.lizbiz.cm » (ancienne version du portail de suivi, encore
en ligne), « test-api.lizbiz.cm » (environnement de test) et
« vpn2.lizbiz.cm » (ancienne passerelle d'accès distant, non retirée).

\medskip
\textbf{Analyse} (confiance moyenne). Les trois serveurs oubliés sont
précisément ceux qui ne reçoivent ni correctifs ni second facteur. Leur
résolution et leur appartenance à la plage du périmètre restent à vérifier
avant tout test actif (Chapitre~\ref{chap:engagement}).
\end{boxscenario}

\begin{boxobjet}{T-03 · Ancienne passerelle d'accès distant non retirée}
\textbf{Description~:} « vpn2.lizbiz.cm » répond sur Internet et accepte une
authentification par mot de passe seul. \textbf{Preuve~:} entrée du journal
de certificats, page de connexion (capture du 21~octobre).
\textbf{Gravité~:} élevée. \textbf{Impact~:} contournement du second facteur
de la passerelle principale. \textbf{Recommandation~:} retirer le service,
révoquer son certificat, inscrire le retrait au processus de fin de vie.
\textbf{Détection~:} inventaire externe périodique.
\end{boxobjet}

% ============================================================================
\section{Les Limites de la Collecte}
\label{sec:c05-limites}
% ============================================================================

\begin{boxalert}
\textbf{Violation typique.} Se connecter à un compte trouvé dans une fuite
publique « pour vérifier qu'il fonctionne ».

\smallskip
\textbf{Règle~:} utiliser un identifiant volé, même publié, est un accès
non autorisé s'il n'est pas prévu par les règles d'engagement. La collecte
passive s'arrête à la lecture~; on signale la fuite au client.
\end{boxalert}

La collecte sur les personnes (salariés, dirigeants) touche à des données
personnelles~: on se limite à ce qui sert le test (format des adresses,
fonctions exposées), on ne constitue pas de dossier individuel, et on
détruit les données à la fin de la mission (Chapitre~\ref{chap:droit}).

% ============================================================================
\section{Réduire l'Exposition}
\label{sec:c05-reduire}
% ============================================================================

\begin{boxprocedure}
\textbf{Mesurer et réduire sa propre exposition}
\begin{enumerate}
    \item Surveiller les journaux de certificats pour ses domaines.
    \item Comparer chaque trimestre la surface externe observée à
          l'inventaire.
    \item Retirer les services anciens et les environnements de test exposés.
    \item Relire les offres d'emploi et les communications techniques
          publiques.
    \item Rechercher ses domaines dans les dépôts de code publics.
\end{enumerate}
\end{boxprocedure}

\begin{boxmetier}{RSSI}
\textbf{Faire faire l'OSINT de sa propre organisation.} Le RSSI commande
chaque trimestre une reconnaissance passive de son organisation~: ce qu'un
attaquant voit gratuitement, il doit le voir avant lui.
\end{boxmetier}

\separateur

\begin{boxresume}
\textbf{Ce qu'il faut retenir de ce chapitre~:}
\begin{itemize}
    \item \textbf{Passive}~: sans toucher la cible~; \textbf{active}~: en
          interagissant.
    \item \textbf{DNS, certificats, registres, annonces, réseaux, dépôts}~:
          les sources ouvertes.
    \item \textbf{La surface externe observée} se compare à l'inventaire~;
          les écarts sont souvent les failles.
    \item \textbf{La collecte s'arrête à la lecture}~; données personnelles
          minimisées.
    \item \textbf{Faire l'OSINT de soi-même}, avant l'attaquant.
\end{itemize}
\end{boxresume}

\begin{boxinfo}
\textbf{Points de contrôle de l'élève.} Je sais~: distinguer reconnaissance
passive et active~; citer six sources ouvertes~; expliquer la transparence
des certificats~; comparer surface observée et inventaire~; respecter les
limites de la collecte.
\end{boxinfo}

\bigskip

\begin{boxexo}[title={\ding{45}\quad Exercice~5.1. Passive ou active \niveaubase}]
Classez~: lire un journal de certificats~; résoudre un nom auprès du
serveur DNS de la cible~; consulter une offre d'emploi~; balayer un port.
\end{boxexo}

\begin{boxexo}[title={\ding{45}\quad Exercice~5.2. L'écart \niveaubase}]
Combien de sous-domaines trouvés ne figuraient pas à l'inventaire de
LiZbiZ~? Quelle part du total~?
\end{boxexo}

\begin{boxexo}[title={\ding{45}\quad Exercice~5.3. Les certificats \niveaubase}]
Pourquoi les journaux de transparence des certificats révèlent-ils des
serveurs « discrets »~?
\end{boxexo}

\begin{boxexo}[title={\ding{45}\quad Exercice~5.4. L'offre d'emploi \niveaumoyen}]
Réécrivez l'offre d'emploi de LiZbiZ pour qu'elle reste attractive sans
révéler ses technologies de sécurité.
\end{boxexo}

\begin{boxexo}[title={\ding{45}\quad Exercice~5.5. Le compte publié \niveaumoyen}]
Vous trouvez dans une fuite publique l'adresse et le mot de passe d'un
salarié de LiZbiZ. Que faites-vous, et que ne faites-vous pas~?
\end{boxexo}

\begin{boxexo}[title={\ding{45}\quad Exercice~5.6. Les personnes \niveaumoyen}]
Quelles informations sur les salariés le testeur peut-il collecter pour une
campagne d'hameçonnage simulée, et lesquelles doit-il s'interdire~?
\end{boxexo}

\begin{boxexo}[title={\ding{45}\quad Exercice~5.7. T-03 \niveaumoyen}]
Expliquez pourquoi l'ancienne passerelle annule l'effet du second facteur
de la passerelle principale.
\end{boxexo}

\begin{boxexo}[title={\ding{45}\quad Exercice~5.8. La procédure \niveaumoyen}]
Transformez la procédure « mesurer et réduire sa propre exposition » en
indicateur clé de risque avec des seuils (cours GES4062).
\end{boxexo}

\begin{boxexo}[title={\ding{45}\quad Exercice~5.9. L'ANPP \niveauexpert}]
Préparez le plan de reconnaissance passive de l'Opération MEFOU~: sources,
informations attendues, limites.
\end{boxexo}

\begin{boxexo}[title={\ding{45}\quad Exercice~5.10. ATT\&CK \niveauexpert}]
Pour la tactique « reconnaissance » d'ATT\&CK, proposez trois mesures
défensives qui la rendent plus difficile ou plus visible.
\end{boxexo}

\bigskip

\begin{boxinfo}
\textbf{Pour aller plus loin.}
\begin{itemize}
    \item RFC~6962~\cite{rfc6962}~; MITRE ATT\&CK~\cite{mitreattack}.
    \item PTES, collecte de renseignements~\cite{ptes}.
\end{itemize}
\end{boxinfo}

\begin{boxinfo}
\textbf{Transition.} La reconnaissance active commence par le balayage.
C'est l'objet du Chapitre~\ref{chap:balayage}.
\end{boxinfo}
